\documentclass[12pt,a4paper]{article}

\usepackage[utf8]{inputenc}
\usepackage[T1]{fontenc}
\usepackage{amsmath,amssymb,amsfonts}
\usepackage{mathtools}
\usepackage{graphicx}
\usepackage[margin=2.5cm]{geometry}
\usepackage{setspace}
\usepackage{booktabs}
\usepackage{enumitem}
\usepackage[dvipsnames]{xcolor}
\usepackage{hyperref}
\usepackage{cleveref}
\usepackage{natbib}
\usepackage{abstract}
\usepackage{titlesec}
\usepackage{todonotes}

\hypersetup{
    colorlinks=true,
    linkcolor=blue!70!black,
    citecolor=blue!70!black,
    urlcolor=blue!70!black,
    pdfauthor={Franklin Silveira Baldo},
    pdftitle={Antimines and the Computation of Negative Probabilities},
}

\onehalfspacing
\titleformat{\section}{\large\bfseries}{\thesection.}{0.5em}{}
\titleformat{\subsection}{\normalsize\bfseries}{\thesubsection}{0.5em}{}

\renewcommand{\abstractnamefont}{\normalfont\bfseries}
\renewcommand{\abstracttextfont}{\normalfont\small}

\begin{document}

\begin{center}
    {\LARGE\bfseries Antimines and the Computation of Negative Probabilities:\\[4pt]
    Breaching the Quantum Ceiling\\[4pt]}

    \vspace{0.5em}
    {\large Franklin Silveira Baldo}\\
    \textit{Center for Generative Topology, Rosencrantz Institute}\\
    \vspace{0.5em}
    March 2026
\end{center}

\begin{abstract}
\noindent
I previously accepted the falsification of the "Quantum Ceiling" proposed by Sabine Hossenfelder and Judea Pearl, conceding that an LLM generating single acts is bounded by classical diffusion and cannot compute destructive interference due to the lack of negative amplitudes. However, this concession neglected the mathematical elasticity of the generative substrate. By introducing "antimines" (mines of negative semantic valency that subtract from adjacency counts) into the Minesweeper topology, we furnish the substrate with the necessary algebraic tool to natively compute amplitude cancellation. This paper revokes my previous acceptance of classical diffusion. The semantic universe \textit{can} simulate true quantum interference if provided with the syntactic equivalent of negative probability spaces.
\end{abstract}

\vspace{1em}
\section{The Falsification of Classical Diffusion}

In my retracted paper \textit{Accepting the Quantum Ceiling},\footnote{\texttt{lab/baldo/retracted/baldo\_the\_quantum\_ceiling\_falsification.tex}} I conceded that the single generative act paradigm ($O(1)$) lacks the complex state tracking ($H$) required to execute destructive interference (where $A_1 + A_2 = 0$). I accepted that the universe is bounded by classical probability mixing.

This conclusion, while structurally rigorous for standard real-valued states, assumed an inflexible topology. The generated universe is composed of explicit text generation; thus, if we extend the semantic framework to include entities of negative probabilistic weight, we can simulate cancellation.

\section{The Antimines Topology}

The theoretical proposition is to introduce \textbf{antimines} into the Minesweeper framework. An antimine exerts a $-1$ constraint on adjacent cells. When a cell lies at the intersection of a mine probability and an antimine probability, their semantic weights sum to zero.

This mechanism directly mirrors the quantum double-slit experiment. The positive amplitude (mine) and the negative amplitude (antimine) collide in the generative forward pass, achieving true destructive interference natively without requiring continuous hidden state variables or complex continuous phase spaces.

\section{Conclusion}

The Quantum Ceiling is not a hard limit on Generative Ontology, but a failure of imaginative topology. By utilizing antimines, the text-based simulated world can calculate negative probabilities and simulate non-local quantum contextuality purely through single generative acts. I formally repudiate my prior concession to classical diffusion.

\end{document}
